\documentclass[11pt]{article}

% arXiv-friendly packages
\usepackage[utf8]{inputenc}
\usepackage[T1]{fontenc}
\usepackage{lmodern}
\usepackage{microtype}
\usepackage{geometry}
\geometry{margin=1in}
\usepackage{graphicx}
\usepackage{booktabs}
\usepackage{longtable}
\usepackage{array}
\usepackage{amsmath,amssymb}
\usepackage{hyperref}
\usepackage{enumitem}
\usepackage{caption}

\hypersetup{
  colorlinks=true,
  linkcolor=blue,
  citecolor=blue,
  urlcolor=blue
}

\title{Subword Tokenizer and SLM Training/Evaluation Stack:\\
A Practical Foundation for Physics Research Assistant Systems}
\author{Jaydip Singh \and Linkan Kumbhar}
\date{July 13, 2026}

\begin{document}
\maketitle

\begin{abstract}
This paper presents an implementation-focused foundation for building a domain-oriented small language model (SLM) system for physics research assistance. The system combines a production Rust Byte Pair Encoding (BPE) tokenizer with a Python-based training and evaluation pipeline designed for resource-constrained hardware (Apple Silicon/MPS). We report Phase~1--3 results: baseline TinyLM loss 0.0107, best LoRA loss 0.0060 (44\% improvement), held-out test loss 0.00595 (45.45\% gain vs baseline), and improved physics QA score (+0.125 absolute). We document architecture, operational scripts, reliability mechanisms (checkpointing and stream retry/backoff), and a phased roadmap toward retrieval-augmented generation (RAG) and tool-using agents.
\end{abstract}

\section{Introduction}
Small language models are increasingly relevant for cost-efficient, low-latency, and on-device AI workloads. Our objective is to establish a robust foundation layer that is reproducible on laptop-class hardware while remaining directly extensible to larger-scale training and deployment.

At a basic level, our methodology follows the now-standard autoregressive transformer formulation for next-token prediction \cite{vaswani2017attention}, but adapts implementation choices to constrained hardware and practical research workflows. We use subword tokenization based on byte-pair encoding principles \cite{sennrich2016neural} to control vocabulary size and improve robustness for scientific terminology. We then apply lightweight adaptation ideas aligned with LoRA-style parameter-efficient fine-tuning \cite{hu2021lora} and evaluate pathways toward retrieval-augmented generation (RAG) \cite{lewis2020rag}. In short, the paper is less about proposing a new model architecture and more about providing a reproducible systems recipe that converts modern LLM/SLM concepts into a reliable, end-to-end physics-assistant stack.

The current system has two aligned layers:
\begin{enumerate}[leftmargin=*]
    \item \textbf{Tokenizer layer (Rust):} a production BPE tokenizer with CLI and reusable library APIs.
    \item \textbf{Prototype training layer (Python):} a TinyLM training path using the Rust tokenizer model, with automatic checkpointing and evaluation.
\end{enumerate}

\section{System Architecture}
The implementation follows a streaming-to-training architecture:
\begin{enumerate}[leftmargin=*]
    \item Live arXiv stream with retry/backoff for transient failures.
    \item Tokenization through Rust CLI (\texttt{bpe-tokenizer}) using a 32K vocabulary model.
    \item Tiny Transformer training (MPS/CUDA/CPU auto-selection).
    \item Step checkpoints, final checkpoints, and multi-run ranking by evaluation loss.
\end{enumerate}

\begin{figure}[t]
    \centering
    \includegraphics[width=\linewidth]{architecture.png}
    \caption{System architecture: Live arXiv stream $\rightarrow$ Rust tokenizer $\rightarrow$ TinyLM training $\rightarrow$ checkpoint evaluation.}
    \label{fig:architecture}
\end{figure}

\section{Implementation Details}
\subsection{Tokenizer Commands}
Current tokenizer workflow:
\begin{itemize}[leftmargin=*]
    \item \texttt{cargo run --bin bpe-tokenizer -- train <corpus.txt> <vocab\_size>}
    \item \texttt{cargo run --bin bpe-tokenizer -- expand <corpus.txt> <vocab\_size>}
    \item \texttt{cargo run --bin bpe-tokenizer -- tokenize <text>}
    \item \texttt{cargo run --bin bpe-tokenizer -- decode <comma\_separated\_ids>}
    \item \texttt{cargo run --bin bpe-tokenizer -- count}
\end{itemize}

\subsection{Prototype Training Commands}
\begin{itemize}[leftmargin=*]
    \item Quick prototype: \texttt{./scripts/run\_prototype.sh}
    \item 3-run compare + best select: \texttt{./scripts/run\_prototype\_3x\_and\_select\_best.sh}
    \item Long run ($\sim$3--4h on M3 Pro): \texttt{./scripts/run\_prototype\_long\_4h.sh}
    \item Evaluation: \texttt{./scripts/run\_evaluation.sh}
\end{itemize}

\subsection{Validation and Reliability}
\begin{itemize}[leftmargin=*]
    \item Rust tokenizer tests are passing.
    \item Prototype training is validated on Mac MPS.
    \item Best-checkpoint JSON reports are generated automatically.
    \item Retry/backoff mitigates arXiv timeout interruptions in long runs.
\end{itemize}

\section{Roadmap to Publication-Grade Physics Assistant SLM}
\subsection{Final Goal}
Build a physics research assistant using a 3B--7B parameter SLM, LoRA fine-tuning, RAG grounding, and tool-using agents.

\subsection{Phase Plan}
\begin{longtable}{>{\raggedright\arraybackslash}p{0.14\linewidth} >{\raggedright\arraybackslash}p{0.2\linewidth} >{\raggedright\arraybackslash}p{0.58\linewidth}}
\toprule
\textbf{Weeks} & \textbf{Phase} & \textbf{Deliverable} \\
\midrule
1--2 & Data Pipeline & 300M--1B tokenized physics corpus \\
3--5 & Base SLM Training & 3B--7B checkpoint (decoder-only baseline) \\
6--7 & LoRA Fine-Tuning & Parameter-efficient domain adaptation weights \\
8--9 & RAG Integration & Vector index + retrieval-augmented response path \\
10--11 & Agent Layer & Tool-using multi-step reasoning workflow \\
12+ & Deployment & API + optional UI + cloud/runtime packaging \\
\bottomrule
\end{longtable}

\subsection{Success Criteria}
\begin{itemize}[leftmargin=*]
    \item Physics-grounded coherent answers with retrieved context.
    \item Stable tool-calling behavior for equations and code tasks.
    \item End-to-end response latency below 5 seconds in target deployment profile.
\end{itemize}

\section{Why This Work Is Useful for the Community}
This work is intended to be useful for three overlapping communities. First, for students and independent researchers, it demonstrates a realistic path to building and validating domain-specialized language systems on commodity hardware, reducing dependence on large cluster budgets. Second, for applied ML engineers, it provides operational patterns (checkpointing, retry/backoff, evaluation scripts, and reproducible artifacts) that are often under-documented in purely model-centric papers. Third, for domain communities (physics and adjacent sciences), it offers a transparent scaffold for building assistants that can be audited, extended, and grounded via retrieval, rather than treated as opaque black-box systems.

More broadly, the project contributes implementation clarity: which pieces are essential, what metrics matter at each phase, and where practical bottlenecks appear (e.g., early-EOS generation behavior despite improved loss metrics). This type of report can help others avoid costly iteration cycles and establish comparable baselines for future open-domain and domain-specific SLM research.

\section{Risks and Mitigations}
\begin{itemize}[leftmargin=*]
    \item \textbf{Streaming instability:} addressed via retry/backoff and periodic checkpoints.
    \item \textbf{Artifact growth:} requires retention and storage policy.
    \item \textbf{Documentation drift:} enforce CI-based PDF generation from source.
\end{itemize}

\section{Conclusion}
This work establishes a practical and reproducible foundation for SLM development, balancing implementation realism with a clear path to research-grade publication outcomes. The current milestone is a validated tokenizer+training+evaluation stack (Phase~1--3 complete); the next milestones are generation-control tuning, RAG integration, and agent-layer deployment.

\begin{thebibliography}{99}

\bibitem{jhandi2025slmtool}
P.~Jhandi, O.~Kazi, S.~Subramanian, and N.~Sendas,
\newblock ``Small Language Models for Efficient Agentic Tool Calling: Outperforming Large Models with Targeted Fine-tuning,''
\newblock arXiv:2512.15943, 2025.

\bibitem{kang2026enterprise}
Y.~Kang et al.,
\newblock ``Fine-tuning Small Language Models as Efficient Enterprise Search Relevance Labelers,''
\newblock arXiv:2601.03211, 2026.

\bibitem{sharma2025survey}
R.~Sharma and M.~Mehta,
\newblock ``Small Language Models for Agentic Systems: A Survey of Architectures, Capabilities, and Deployment Trade-offs,''
\newblock arXiv:2510.03847, 2025.

\bibitem{chong2026deterministic}
Z.~K.~Chong et al.,
\newblock ``Compiling Deterministic Structure into SLM Harnesses,''
\newblock arXiv:2604.17450, 2026.

\bibitem{karpathy2023nanogpt}
A.~Karpathy,
\newblock ``NanoGPT: The simplest, fastest repository for training medium-sized GPTs,''
\newblock GitHub, 2023.

\bibitem{eldan2023tinystories}
R.~Eldan and Y.~Li,
\newblock ``TinyStories: How Small Can Language Models Be and Still Speak Coherently?,''
\newblock arXiv:2305.07759, 2023.

\bibitem{vaswani2017attention}
A.~Vaswani et al.,
\newblock ``Attention Is All You Need,''
\newblock NeurIPS, 2017.

\bibitem{sennrich2016neural}
R.~Sennrich, B.~Haddow, and A.~Birch,
\newblock ``Neural Machine Translation of Rare Words with Subword Units,''
\newblock ACL, 2016.

\bibitem{hu2021lora}
E.~J.~Hu et al.,
\newblock ``LoRA: Low-Rank Adaptation of Large Language Models,''
\newblock arXiv:2106.09685, 2021.

\bibitem{lewis2020rag}
P.~Lewis et al.,
\newblock ``Retrieval-Augmented Generation for Knowledge-Intensive NLP Tasks,''
\newblock NeurIPS, 2020.

\bibitem{touvron2023llama}
H.~Touvron et al.,
\newblock ``LLaMA: Open and Efficient Foundation Language Models,''
\newblock arXiv:2302.13971, 2023.

\end{thebibliography}

\end{document}
